% GENERATED FILE. Edit canonical lessons, then run npm run generate:books.
% canonical-source-hash: fnv1a64:7633c08c42528f89
% canonical-lessons: UR-C69-hathi, UR-C69-bandar, UR-C69-sap, UR-C69-richh, UR-C69-bara

\chapter{Elephant, Monkey, Snake, Bear, Big}
\label{ch:ur-elephant-monkey-snake-bear-big}
\hlchaptermodality{69}

\begin{chapteropening}
\textbf{By the end of this chapter:} \emph{I can say an elephant, a monkey, a snake, a bear and big.}

Complete the last lesson of chapter 69: I can say an elephant, a monkey, a snake, a bear and big.
\end{chapteropening}

\section[\texorpdfstring{\ur{ہاتھی} (hāthī)}{ہاتھی (hāthī)}]{\ur{ہاتھی} (hāthī) — an elephant}

\label{lesson:UR-C69-hathi}



\begin{quote}
Before the new one: say the Urdu for a mouse, then the Urdu for a lion.
\end{quote}

\subsection*{You'll want to know: \ur{ہاتھی}}
\textbf{\ur{ہاتھی}} — \emph{hāthī} — \textquotedblleft{}an elephant\textquotedblright{}.

\begin{cousinweb}
\textbf{\ur{ہاتھی}} (\emph{hāthī}) is a tadbhava, a word inherited from Sanskrit and reshaped by sound change. It comes from \emph{hastin}, \textquotedblleft{}having hands; elephant\textquotedblright{} (CDIAL 14039). The elephant was \textquotedblleft{}the one with a hand\textquotedblright{}, its trunk.
\end{cousinweb}

\begin{grammarlens}[title={What you've built}]
One more word for this chapter.
\end{grammarlens}

\subsection*{Your turn}
\begin{itemize}
\raggedright
  \item \emph{Say it:} \emph{hāthī}
  \item \emph{Say it:} \emph{hāthī}, once more
  \item \emph{Say it:} say \emph{hāthī}
  \item \emph{Recall:} say the Urdu for a mouse, then the Urdu for a lion, then say \emph{hāthī} again
\end{itemize}

\begin{tcolorbox}[breakable,colback=teal!4,colframe=teal!35!black,title={Before you move on}]
What does \ur{ہاتھی} mean? (\textquotedblleft{}An elephant\textquotedblright{}.) Say it once more.
\end{tcolorbox}
\section[\texorpdfstring{\ur{بندر} (bandar)}{بندر (bandar)}]{\ur{بندر} (bandar) — a monkey}

\label{lesson:UR-C69-bandar}



\begin{quote}
Before the new one: say the Urdu for a lion, then the Urdu for an elephant.
\end{quote}

\subsection*{You'll want to know: \ur{بندر}}
\textbf{\ur{بندر}} — \emph{bandar} — \textquotedblleft{}a monkey\textquotedblright{}.

\begin{cousinweb}
\textbf{\ur{بندر}} (\emph{bandar}) is a tadbhava, a word inherited from Sanskrit and reshaped by sound change. It comes from \emph{vānara}, \textquotedblleft{}monkey\textquotedblright{} (CDIAL 11515). CDIAL derives the Sanskrit word from \emph{vana-}, \textquotedblleft{}forest\textquotedblright{}.
\end{cousinweb}

\begin{grammarlens}[title={What you've built}]
One more word for this chapter.
\end{grammarlens}

\subsection*{Your turn}
\begin{itemize}
\raggedright
  \item \emph{Say it:} \emph{bandar}
  \item \emph{Say it:} \emph{bandar}, once more
  \item \emph{Say it:} say \emph{bandar}
  \item \emph{Recall:} say the Urdu for a lion, then the Urdu for an elephant, then say \emph{bandar} again
\end{itemize}

\begin{tcolorbox}[breakable,colback=teal!4,colframe=teal!35!black,title={Before you move on}]
What does \ur{بندر} mean? (\textquotedblleft{}A monkey\textquotedblright{}.) Say it once more.
\end{tcolorbox}
\section[\texorpdfstring{\ur{سانپ} (sā̃p)}{سانپ (sā̃p)}]{\ur{سانپ} (sā̃p) — a snake}

\label{lesson:UR-C69-sap}



\begin{quote}
Before the new one: say the Urdu for an elephant, then the Urdu for a monkey.
\end{quote}

\subsection*{You'll want to know: \ur{سانپ}}
\textbf{\ur{سانپ}} — \emph{sā̃p} — \textquotedblleft{}a snake\textquotedblright{}.

\begin{cousinweb}
\textbf{\ur{سانپ}} (\emph{sā̃p}) is a tadbhava, a word inherited from Sanskrit and reshaped by sound change. It comes from \emph{sarpa}, \textquotedblleft{}snake\textquotedblright{} (CDIAL 13271).
\end{cousinweb}

\begin{grammarlens}[title={What you've built}]
One more word for this chapter.
\end{grammarlens}

\subsection*{Your turn}
\begin{itemize}
\raggedright
  \item \emph{Say it:} \emph{sā̃p}
  \item \emph{Say it:} \emph{sā̃p}, once more
  \item \emph{Say it:} say \emph{sā̃p}
  \item \emph{Recall:} say the Urdu for an elephant, then the Urdu for a monkey, then say \emph{sā̃p} again
\end{itemize}

\begin{tcolorbox}[breakable,colback=teal!4,colframe=teal!35!black,title={Before you move on}]
What does \ur{سانپ} mean? (\textquotedblleft{}A snake\textquotedblright{}.) Say it once more.
\end{tcolorbox}
\section[\texorpdfstring{\ur{ریچھ} (rīchh)}{ریچھ (rīchh)}]{\ur{ریچھ} (rīchh) — a bear}

\label{lesson:UR-C69-richh}



\begin{quote}
Before the new one: say the Urdu for a monkey, then the Urdu for a snake.
\end{quote}

\subsection*{You'll want to know: \ur{ریچھ}}
\textbf{\ur{ریچھ}} — \emph{rīchh} — \textquotedblleft{}a bear\textquotedblright{}.

\begin{cousinweb}
\textbf{\ur{ریچھ}} (\emph{rīchh}) is a tadbhava, a word inherited from Sanskrit and reshaped by sound change. It comes from \emph{ṛkṣa}, \textquotedblleft{}bear\textquotedblright{} (CDIAL 2445).
\end{cousinweb}

\begin{grammarlens}[title={What you've built}]
One more word for this chapter.
\end{grammarlens}

\subsection*{Your turn}
\begin{itemize}
\raggedright
  \item \emph{Say it:} \emph{rīchh}
  \item \emph{Say it:} \emph{rīchh}, once more
  \item \emph{Say it:} say \emph{rīchh}
  \item \emph{Recall:} say the Urdu for a monkey, then the Urdu for a snake, then say \emph{rīchh} again
\end{itemize}

\begin{tcolorbox}[breakable,colback=teal!4,colframe=teal!35!black,title={Before you move on}]
What does \ur{ریچھ} mean? (\textquotedblleft{}A bear\textquotedblright{}.) Say it once more.
\end{tcolorbox}
\section[\texorpdfstring{\ur{بڑا} (baṛā)}{بڑا (baṛā)}]{\ur{بڑا} (baṛā) — big}

\label{lesson:UR-C69-bara}



\begin{quote}
Before the new one: say the Urdu for a snake, then the Urdu for a bear.
\end{quote}

\subsection*{You'll want to know: \ur{بڑا}}
\textbf{\ur{بڑا}} — \emph{baṛā} — \textquotedblleft{}big\textquotedblright{}.

\begin{cousinweb}
\textbf{\ur{بڑا}} (\emph{baṛā}) is a tadbhava, a word inherited from Sanskrit and reshaped by sound change. It comes from \emph{vaḍra}, \textquotedblleft{}big\textquotedblright{} (CDIAL 11225). CDIAL notes that the Sanskrit word is known only from the old word-lists, and doubts a link with \emph{vṛddha}, \textquotedblleft{}grown\textquotedblright{}.
\end{cousinweb}

\begin{grammarlens}[title={What you've built}]
One more word for this chapter.
\end{grammarlens}

\subsection*{Your turn}
\begin{itemize}
\raggedright
  \item \emph{Say it:} \emph{baṛā}
  \item \emph{Say it:} \emph{baṛā}, once more
  \item \emph{Say it:} say \emph{baṛā}
  \item \emph{Recall:} say the Urdu for a snake, then the Urdu for a bear, then say \emph{baṛā} again
\end{itemize}

\begin{tcolorbox}[breakable,colback=teal!4,colframe=teal!35!black,title={Before you move on}]
What does \ur{بڑا} mean? (\textquotedblleft{}Big\textquotedblright{}.) Say it once more.
\end{tcolorbox}
